%% ============================================================
%%  datos.tex  —  Sección 1.3
%% ============================================================

\subsection{Datos utilizados}

\subsubsection{Origen de los datos}

Todos los datos utilizados en este proyecto son \textbf{sintéticos}, generados
proceduralmente mediante la función \texttt{generateRandomGraph(n, m, maxW,
seed)} definida en \texttt{graph.h}. No se utilizaron conjuntos de datos
externos, repositorios públicos ni grafos reales.

La elección de datos sintéticos responde a varios factores:

\begin{itemize}
  \item \textbf{Control experimental:} permite variar de forma independiente
    y exacta los parámetros $n$, $m$ y $w_{\max}$.
  \item \textbf{Reproducibilidad determinista:} la semilla del generador
    Mersenne Twister (\texttt{std::mt19937}) garantiza grafos idénticos entre
    ejecuciones.
  \item \textbf{Escalabilidad:} es posible generar instancias de cualquier
    tamaño sin restricciones de disponibilidad o formato.
\end{itemize}

\subsubsection{Descripción del conjunto de datos}

Cada grafo generado es un \textbf{grafo dirigido} con pesos enteros en el
rango $[1, w_{\max}]$ donde $w_{\max} = 1\,000$ en todos los experimentos.
El algoritmo de generación garantiza la \textit{conectividad débil} del grafo
mediante los siguientes pasos:

\begin{enumerate}
  \item Se genera una permutación aleatoria de los $n$ vértices.
  \item Se insertan $n-1$ aristas dirigidas entre vértices consecutivos de
    dicha permutación, formando un árbol de expansión dirigido.
  \item Se agregan aristas adicionales de forma aleatoria (con auto-bucles
    descartados) hasta alcanzar el objetivo \texttt{targetEdges}, con un
    límite de $10 \times \texttt{targetEdges}$ intentos para evitar bucles
    infinitos en grafos muy pequeños.
\end{enumerate}

\noindent El número de aristas real $m_{\text{real}}$ puede diferir
ligeramente del objetivo $m_{\text{target}}$ si el número de intentos se
agota (lo que ocurre raramente en los rangos evaluados).

La Tabla~\ref{tab:rangos} resume los rangos de los parámetros empleados en
el proyecto.

\begin{table}[H]
  \centering
  \caption{Rango de parámetros de los grafos generados.}
  \label{tab:rangos}
  \begin{tabular}{lll}
    \toprule
    \textbf{Parámetro} & \textbf{Rango} & \textbf{Experimentos} \\
    \midrule
    Vértices $n$           & $100$ a $10\,000$          & Exps.\ 1, 2, 3 \\
    Aristas $m$            & $200$ a $200\,000$         & Exps.\ 1, 2, 3 \\
    Densidad $m/(n(n-1))$  & $0.02\%$ a $20\%$         & Exp.\ 3 \\
    Peso máximo $w_{\max}$ & $1\,000$                   & Todos \\
    Semilla base           & $42$                       & Exps.\ 1, 2, 3 \\
    \midrule
    Vértices $n$ (densidad)  & $500$ a $5\,000$         & Exp.\ densidad \\
    Semillas (densidad)      & $42$, $123$, $999$       & Exp.\ densidad \\
    Densidades (densidad)    & $m \in \{n,\,2n,\,5n,\,10n,\,\lfloor n^2/4\rfloor\}$
                             & Exp.\ densidad \\
    \bottomrule
  \end{tabular}
\end{table}

\subsubsection{Preprocesamiento de datos}

No se aplica ningún preprocesamiento sobre los datos de entrada, ya que los
grafos se generan directamente en la estructura de listas de adyacencia
(\texttt{Graph::adj}) en el formato requerido por los algoritmos.

Los datos de \textit{salida} (tiempos y operaciones) se exportan a archivos
CSV sin transformación adicional. El archivo \texttt{resultados/densidad\_raw.csv}
contiene las 60 mediciones individuales ($4$ valores de $n$ $\times$
$5$ densidades $\times$ $3$ semillas), y
\texttt{resultados/densidad\_promedio.csv} contiene los promedios y
desviaciones estándar calculados directamente en C++ durante la ejecución del
benchmark (\texttt{density\_experiment.cpp}).
